\chapter{Textures, Render Targets and Surfaces}
\label{ch:textures-rendertargets}

This chapter is the one in Part~III where the gap between ``the API exists'' and ``the API
does what it says on every backend'' is widest. \cnaclass{Texture2D} is close to fully
faithful; \cnaclass{Texture3D} and \cnaclass{TextureCube} carried a real, load-bearing
structural deviation from FNA for a stretch of this project's history, closed since (see
\S\ref{sec:texture3d-inheritance}) but worth the full account of what changed and what did
not; and render targets --- once the sharpest per-backend divergence in the whole Graphics
namespace, before a wave of later fixes closed most of it, covered in full below --- still
carry two genuinely open gaps of their own. All three are covered here together, with the
specific per-backend facts intact rather than smoothed into ``mostly works.''

\section{Texture: the base, and one format that actually works}

\cnaclass{Texture} (inheriting \cnaclass{GraphicsResource}) exposes \texttt{Format}
(\cnaclass{SurfaceFormat}) and \texttt{LevelCount}, plus a set of static helpers ported from
FNA's own internal size-calculation methods --- made \texttt{public static} in CNA rather
than FNA's \texttt{internal}, because three of the four real call sites in CNA are not even
\cnaclass{Texture} subclasses. The single most consequential method here is
\texttt{ValidateFormat(format)}: it throws for anything other than
\cnaclass{SurfaceFormat::Color}. This is enforced at construction time, across
\cnaclass{Texture2D}, \cnaclass{Texture3D}, and \cnaclass{TextureCube} alike --- meaning the
other 26 \cnaclass{SurfaceFormat} values (the original 20 XNA values plus 7 FNA-only
extensions like \texttt{ColorBgraEXT}, \texttt{ColorSrgbEXT}, and the block-compressed
\texttt{Bc7EXT} family) cannot be used to construct a texture at all today, regardless of
backend. \texttt{Dxt1} / \texttt{Dxt3} / \texttt{Dxt5} are the one partial exception, but only
through the \texttt{.dds} loading path in \texttt{Texture2D::FromStream}
(\S\ref{sec:texture-streams}), which CPU-decompresses to RGBA8 rather than uploading a real
compressed texture --- an eight-times VRAM overhead, and no native GPU compressed path exists
anywhere in the project.

\section{A structural deviation that closed since this chapter's own first pass}
\label{sec:texture3d-inheritance}

Real XNA's \texttt{Texture3D} and \texttt{TextureCube} both inherit \texttt{Texture}, which
is what lets any of them sit in \texttt{GraphicsDevice.Textures[slot]} and be bound to a
custom shader like an ordinary 2D texture. An earlier pass of this chapter reported that
CNA's \cnaclass{Texture3D} and \cnaclass{TextureCube} did \textbf{not} inherit
\cnaclass{Texture} at all, tracked as Task~863 --- reading the real headers directly today
shows that claim is stale: \texttt{Texture3D.hpp} and \texttt{TextureCube.hpp} both now
declare \texttt{class Texture3D : public Texture} / \texttt{class TextureCube : public
Texture}, matching real XNA exactly. \texttt{Texture3DTests.cpp}'s own regression coverage
names precisely what this fix restores --- assignment into a \cnaclass{TextureCollection}
(\texttt{GraphicsDevice.Textures[slot] = texture3D}), a real compile-time impossibility
before Task~863 closed, plus a correctness follow-on: \texttt{Texture3D::Dispose(bool)} now
calls through to \texttt{Texture::Dispose(bool)} after releasing its own backend handle, so
disposing a bound \cnaclass{Texture3D} correctly unbinds it from
\texttt{GraphicsDevice.Textures} / \texttt{VertexTextures}, the same unbind-on-dispose contract
\cnaclass{Texture2D} already honored and \cnaclass{Texture3D} previously skipped entirely.

The practical consequence this earlier pass described --- ``a custom
\cnaclass{ShaderEffect} has no way to bind a \cnaclass{Texture3D} or \cnaclass{TextureCube}''
--- is also closed, and by a more direct route than the generic texture-slot mechanism the
class-hierarchy fix alone would suggest: \cnaclass{ShaderEffect} gained dedicated
\texttt{SetTexture(int, TextureCube\&)} (Task~1081) and \texttt{SetTexture(int, Texture3D\&)}
(Task~863) overloads, each binding straight to the matching GL texture target
(\texttt{glBindTexture(GL\_TEXTURE\_CUBE\_MAP, ...)} / \texttt{GL\_TEXTURE\_3D}) rather than
routing through the ordinary 2D texture-slot path. This is EasyGL-only as of this writing ---
the other nine backends were explicitly deferred, per the same task's own scope note --- so
``a custom shader can sample a volume or cube texture'' is currently a real, tested EasyGL
capability, not yet a cross-backend one. \cnaclass{EnvironmentMapEffect}'s own
\texttt{EnvironmentMap} field, mentioned in an earlier pass as the \emph{only} way any effect
could reach a cube texture, is now one of two paths rather than the sole exception.

\section{Texture2D}

\cnaclass{Texture2D} is copyable and movable (its backend is a \texttt{shared\_ptr}-held
\texttt{ITextureBackend}), constructible from raw dimensions, from a file path (two
\texttt{NOXNA} overloads --- real XNA has no such constructor, since real XNA loads textures
exclusively through the compiled content pipeline), or via \texttt{FromStream} (device plus
stream, or a five-argument form that resizes/crops to an explicit width, height, and
\texttt{zoom} fit-vs-cover flag). \texttt{SetData} / \texttt{GetData} cover full-array and
partial (level, rectangle, start index, element count) forms, all \cnaclass{Color}-only.
\texttt{SaveAsPng} / \texttt{SaveAsJpeg} round out the surface, alongside several
\texttt{NOXNA} escape hatches (\texttt{GetBackend()}, \texttt{GetCpuPixelsWeak()},
\texttt{CreateFromPixels}) that exist for interop with CNA's own content-pipeline and testing
code rather than for game code.

\subsection{A worked example: building a texture from raw pixel data}

Every one of the API traits described above is visible in one small, real pattern: building a
procedural texture entirely from \cnaclass{Color} data, with no file on disk at all --- a
checkerboard, grounded directly in the real constructor and \texttt{SetData} / \texttt{GetData}
signatures in \texttt{Texture2D.hpp}:

\begin{lstlisting}
const int size = 64;
Texture2D checker(getGraphicsDeviceProperty(), size, size); // defaults to SurfaceFormat::Color

std::vector<Color> pixels(size * size);
for (int y = 0; y < size; ++y) {
    for (int x = 0; x < size; ++x) {
        bool light = ((x / 8) + (y / 8)) % 2 == 0;
        pixels[y * size + x] = light ? Color::White : Color::Black;
    }
}
checker.SetData(pixels.data(), static_cast<int>(pixels.size()));

// Read a single pixel back to verify -- GetData is a real round trip, not a cache echo:
Color topLeft;
checker.GetData(&topLeft, 0, 1);
\end{lstlisting}

Two constraints from this chapter's own opening section apply directly to this example, not
just in the abstract: the two-argument \texttt{Texture2D(device, width, height)} constructor
defaults to \cnaclass{SurfaceFormat::Color}, the one format \texttt{ValidateFormat} actually
accepts --- passing any of the other 26 \cnaclass{SurfaceFormat} values to the explicit-format
constructor overload throws immediately, before a single pixel is ever uploaded. And because
\texttt{SetData} / \texttt{GetData} are \cnaclass{Color}-only (no generic-template overload for
an arbitrary pixel struct, unlike real XNA's \texttt{T[]} generic form), a game working with
any other in-memory pixel representation must convert to \cnaclass{Color} before calling
\texttt{SetData} --- there is no faster typed path around it.

\subsection{FromStream format support}
\label{sec:texture-streams}

Format detection in \texttt{FromStream} starts with a minimal internal DDS-header sniffer
(recognizing only the DXT1/3/5 FourCC codes, decoding via an internal \texttt{DxtUtil}), and
falls through to SDL3\_image's \texttt{IMG\_Load\_IO} --- which auto-detects the container
format from the byte stream itself --- for everything else. Verified round-trip formats:
PNG, JPEG (lossy, within tolerance), 24-bit uncompressed BMP, and DDS (DXT1/3/5). AVIF, TIFF,
and WebP are present in the linked SDL3\_image build's own dependencies but are not
independently tested by CNA's own suite. The \texttt{zoom} parameter on the five-argument
overload matters more than its name suggests: \texttt{zoom=false} uses a simplified
largest-dimension-fit heuristic (not a general \texttt{min(w/w0, h/h0)} bounding-box fit --- it
assumes a roughly square target box), while \texttt{zoom=true} scales up and center-crops to
exactly fill the requested box.

\section{Texture3D and TextureCube}

Beyond the structural deviation described above, \cnaclass{Texture3D} and
\cnaclass{TextureCube} are both non-copyable (each owns a unique GPU handle) but movable ---
\cnaclass{Texture3D} specifically gained move construction/assignment to satisfy a glTF
content-pipeline reader's needs. Neither type's requested \cnaclass{SurfaceFormat} is honored
by any backend; both are always stored as RGBA8 on the GPU regardless of what was requested.

\texttt{docs/texture3d-texturecube-support.md} is candid about the deepest gap here:
\textbf{no backend implements a CPU shadow buffer for \texttt{GetData} readback} on either
type, the way \cnaclass{Texture2D} does. \textbf{EasyGL} genuinely reads pixels back, via a
real per-slice or per-face framebuffer plus \texttt{glReadPixels}. \textbf{Vulkan} and
\textbf{BGFX} both make \texttt{GetData} a total, silent no-op --- the caller's buffer is
left completely untouched, with no exception and no error of any kind, because the call falls
through to an unimplemented base-class default. This escaped detection for a long time
because the existing test suites for these methods only ever asserted that the call did not
crash, never that the returned data was actually correct.

Mip levels above~0 have the identical bug shape on every cell of the matrix except one:
\textbf{EasyGL's \texttt{TextureCube}} alone pre-allocates real per-face storage per mip
level (a fix from an earlier task); \textbf{EasyGL's \texttt{Texture3D}} and \textbf{both
types on Vulkan and BGFX} all hardcode a single mip level regardless of what is requested.
\texttt{CubeMapFace} validation, by contrast, is a place CNA is \emph{stricter} than FNA: an
out-of-range face throws \texttt{std::out\_of\_range} at the public API layer, a safety check
real XNA never performs at all.

\subsection{A worked example: six faces, six SetData calls}

\cnaclass{TextureCube}'s constructor takes a single face size (every face is square and
identical in size, per the real XNA contract) plus \texttt{mipMap} and \texttt{format} --- and,
per this chapter's own opening finding, \texttt{format} must be
\cnaclass{SurfaceFormat::Color} or construction throws immediately. Populating a solid-color
test cube (a stand-in for a real skybox's six photographic faces) is a direct, mechanical loop
over \cnaclass{CubeMapFace}'s six enumerators, since \texttt{SetData} takes the face as its
first argument:

\begin{lstlisting}
const int faceSize = 128;
TextureCube skybox(getGraphicsDeviceProperty(), faceSize, /*mipMap=*/false,
                    SurfaceFormat::Color);

const CubeMapFace faces[6] = {
    CubeMapFace::PositiveX, CubeMapFace::NegativeX,
    CubeMapFace::PositiveY, CubeMapFace::NegativeY,
    CubeMapFace::PositiveZ, CubeMapFace::NegativeZ,
};
const Color faceColors[6] = {
    Color::Red,    Color(128, 0, 0, 255),   // +X, -X
    Color::Green,  Color(0, 128, 0, 255),   // +Y, -Y
    Color::Blue,   Color(0, 0, 128, 255),   // +Z, -Z
};

std::vector<Color> pixels(faceSize * faceSize);
for (int i = 0; i < 6; ++i) {
    std::fill(pixels.begin(), pixels.end(), faceColors[i]);
    skybox.SetData(faces[i], pixels.data(), static_cast<int>(pixels.size()));
}
\end{lstlisting}

One of this chapter's own already-stated facts constrains what this loop can safely do next.
\texttt{GetData} on any of these six faces is a genuine no-op on Vulkan and BGFX (silently
leaves the caller's buffer untouched) --- so a readback-based unit test of this exact loop can
only run meaningfully on EasyGL, which is exactly why \texttt{docs/texture3d-texturecube-support.md}
scopes its own \texttt{GetData} verification to that one backend. \texttt{skybox} \emph{can}
now be bound to a custom \cnaclass{ShaderEffect} on EasyGL, per
\S\ref{sec:texture3d-inheritance} above --- either through \texttt{SetTexture(int,
TextureCube\&)} directly, or, since \cnaclass{TextureCube} now inherits \cnaclass{Texture},
through the ordinary \texttt{GraphicsDevice.Textures[slot]} assignment every 2D texture in
this chapter already uses.

\subsection{A worked example: sampling a volume texture from a custom shader}
\label{sec:texture3d-shader-example}

\texttt{examples/easygl\_shadereffect\_texture3d\_test.cpp} is the real, EasyGL-only test
that proves \texttt{ShaderEffect::SetTexture(int, Texture3D\&)} genuinely reaches a
\texttt{sampler3D} uniform, not just that the call compiles. The volume is deliberately
tiny --- a $1\times1\times2$ texture, one texel per Z-slice --- so the whole capability proof
reduces to two texel reads instead of a real ray-marching setup:

\begin{lstlisting}
Texture3D volumeTex(device, 1, 1, 2, /*mipMap=*/false, SurfaceFormat::Color);
const Color slices[2] = {
    Color(255, 0, 0, 255), // Z=0
    Color(0, 0, 255, 255), // Z=1
};
volumeTex.SetData(slices, 2);

auto* fx = dynamic_cast<ShaderEffect*>(fxBase.get());
fx->Apply();
fx->SetTexture(0, volumeTex);
fx->SetUniformInt("VolumeSampler", 0);
fx->SetUniformVec3("coord", 0.5f, 0.5f, 0.25f); // samples Z-slice 0's texel centre
\end{lstlisting}

Sampling at $Z=0.25$ and $Z=0.75$ --- each slice's own texel centre, per
$(\text{sliceIndex} + 0.5) / \text{depth}$ --- reads back \texttt{(255,0,0,255)} and
\texttt{(0,0,255,255)} respectively, even with GL\_LINEAR filtering on both axes: linear
interpolation sampled exactly at a texel centre puts its full weight on that one texel, so
there is no blending artifact to account for. The real test's own false-positive guard is
worth naming as a technique in its own right, since it recurs throughout this book's worked
examples: a second, all-black ``decoy'' \cnaclass{Texture3D} is created and uploaded
\emph{after} the real one, so an unmutated \texttt{SetTexture()} call would leave GL's
texture-unit-0/\texttt{GL\_TEXTURE\_3D} binding pointing at the decoy rather than the
intended volume --- the two-slice colour check only passes for the right reason (a genuine,
per-call rebind) if disabling the \texttt{SetTexture()} call would make it fail.

\subsection{A worked example: rendering into one RenderTargetCube face}

\cnaclass{GraphicsDevice} has a dedicated overload for exactly this purpose ---
\texttt{SetRenderTarget(RenderTargetCube*, CubeMapFace)} --- distinct from the plain
\texttt{SetRenderTarget(RenderTarget2D*)} form used for ordinary 2D targets, and from the
\texttt{SetRenderTargets(const std::vector<RenderTargetBinding>\&)} form used for
simultaneous multi-target binding (\cnaclass{RenderTargetBinding}'s own
\texttt{(Texture*, CubeMapFace)} constructor exists specifically to build entries for that
vector form, not for single-target binding). Rendering a small scene into a single cube face
and then reading it back is the pattern a real-time reflection probe or a baked-environment-map
tool would use:

\begin{lstlisting}
RenderTargetCube probe(getGraphicsDeviceProperty(), 256, /*mipMap=*/false,
                        SurfaceFormat::Color, DepthFormat::Depth24Stencil8);

device.SetRenderTarget(&probe, CubeMapFace::PositiveZ);
device.Clear(Color::CornflowerBlue);
// ... draw the scene as seen looking down +Z from the probe's location ...
device.SetRenderTarget(nullptr); // back to the backbuffer
\end{lstlisting}

Whether this actually produces a correct image is, per \S\ref{sec:rendertargets} below,
backend-dependent in a way no amount of correct call-site code can work around: EasyGL renders
this correctly; Vulkan renders black for either (or both) of two distinct, separately-tracked
reasons; and BGFX has an unconditional \texttt{static\_cast} type-confusion bug reading the
wrong handle pool entirely. A porting engineer hitting a black cube face on Vulkan should treat
this as a confirmed, known backend limitation rather than a bug in call-site code structured
like the example above.

\section{Render targets: a once-sharp per-backend divergence, mostly closed since}
\label{sec:rendertargets}

\cnaclass{RenderTarget2D} inherits both \cnaclass{Texture2D} and \cnaclass{IRenderTarget},
which is what lets it be both drawn into and sampled as an ordinary texture afterward; its
\texttt{Dispose(bool)} correctly throws \texttt{InvalidOperationException} if it is still
bound to the device, matching FNA. \cnaclass{RenderTargetCube}, by contrast, inherits
\cnaclass{GraphicsResource} directly rather than \cnaclass{Texture2D} --- a downstream
consequence of the \cnaclass{TextureCube} structural deviation above --- and so it
\emph{cannot} get the same bound-while-disposing guard, since \cnaclass{RenderTargetBinding}
only ever stores a \texttt{Texture*}. This is tracked, not fixed.

\texttt{docs/rendertarget-support.md} once listed five further per-backend divergences here
(Tasks~877--881) as open gaps, and an earlier pass of this chapter reproduced that list
faithfully. Reading every backend's own current \texttt{.cpp} directly instead of trusting that
document's status column finds four of the five already closed --- each superseded by a later
task (\texttt{Task~903} / \texttt{907} / \texttt{911}) whose own fix comments the tracking document
was never updated to reflect, the identical ``a status that looked settled turned out to be
stale'' pattern Chapters~\ref{ch:easygl-backend} and~\ref{ch:vulkan-backend} both ran into
independently while auditing this exact same document for their own backends. What is actually
still open, and what is now fixed:

\begin{description}
  \item[Sampling a \cnaclass{RenderTargetCube} face after rendering into it, then unbinding.]
        \textbf{EasyGL}: works correctly. \textbf{Vulkan}: still renders black, for the two
        distinct, separately-tracked reasons Chapter~\ref{ch:vulkan-backend} covers in full ---
        a clear-only render target never actually records a render pass, and a second, separate
        black-render defect whose root cause remains genuinely unresolved between two named
        candidates. \textbf{BGFX}: fixed --- the unconditional \texttt{static\_cast} that once
        read a framebuffer-pool handle where a texture-pool handle was expected
        (Chapter~\ref{ch:bgfx-backend} names the exact mechanism) is gone, replaced by a real
        \texttt{dynamic\_cast} to a shared cube-samplable interface for the
        \cnaclass{EnvironmentMapEffect} path and virtual width/height dispatch for the ordinary
        \cnaclass{SpriteBatch} path.
  \item[Depth/stencil format fidelity.] Now real on all three backends: each maps the actual
        requested \cnaclass{DepthFormat} ordinal to the correct native format (GL internal
        format plus attachment point on EasyGL; a per-instance Vulkan format; bgfx's own
        format enum), rather than the single, always-the-same format every backend once
        allocated regardless of what was asked for.
  \item[Mipmap generation.] \textbf{EasyGL} already worked --- real per-level GPU storage,
        auto-regenerated via \texttt{glGenerateMipmap} on unbind, matching FNA3D's own real
        mechanism. \textbf{Vulkan} and \textbf{BGFX} now do too: a real, per-level mip-chain
        regeneration cascade (\texttt{vkCmdBlitImage} on Vulkan; bgfx's own equivalent) runs
        against a render target's just-rendered content once its render pass ends, rather than
        silently accepting and discarding the request the way both once did.
  \item[Multisample antialiasing.] \textbf{EasyGL} already worked: a real multisampled
        renderbuffer with a \texttt{glBlitFramebuffer} resolve on unbind. \textbf{Vulkan} and
        \textbf{BGFX} now genuinely implement it too, rather than honestly reporting
        \texttt{MultiSampleCount = 0} the way both once correctly did while the feature was
        still unimplemented.
\end{description}

Two further, project-wide findings from the same original list are also now closed, not open.
Switching the active render target correctly resets \texttt{Viewport} and
\texttt{ScissorRectangle} to the new target's (or the backbuffer's) own dimensions, on every
backend, and \texttt{GraphicsDevice.Viewport} itself now has real GPU wiring on all three
hardware backends investigated --- EasyGL applies a custom sub-region \cnaclass{Viewport}
unconditionally, whether or not a render target is currently bound, while Vulkan's own version
of the identical fix is honestly narrower, backbuffer-pass-only, for an architectural reason
Chapters~\ref{ch:vulkan-backend} and~\ref{ch:easygl-backend} both cover in full. And
\texttt{GraphicsDevice::SetRenderTargets()} now genuinely enforces FNA's real four-target
\texttt{MAX\_RENDERTARGET\_BINDINGS} cap at the shared, cross-backend level --- throwing if a
caller passes more than four bindings, regardless of what any individual backend's own internal
buffer happens to be sized for --- closing what this section once described as no backend
matching that real XNA limit at all.

\section{SurfaceFormat, in one sentence}

\texttt{docs/surface-format-support.md} boils the entire 27-value \cnaclass{SurfaceFormat}
enum down to a single, load-bearing fact: only \cnaclass{SurfaceFormat::Color} is genuinely
supported by any backend today, everything else is silently treated as RGBA8 wherever it is
not rejected outright by \texttt{ValidateFormat} at construction time. One real, historical
bug is worth naming because of what it reveals about how subtle a graphics-correctness bug
can be: Vulkan's texture path once created images as \texttt{VK\_FORMAT\_R8G8B8A8\_SRGB}
(wrong --- \cnaclass{SurfaceFormat::Color} is linear, not sRGB) while the swapchain
\emph{separately} preferred a matching sRGB format. The two wrong encode/decode steps
approximately canceled out for ordinary textured content, which is exactly why the bug went
unnoticed for a long time --- it only became visible on non-textured content (flat vertex
colors, raw lighting output), where a mid-gray value of 128 was read back as 188. Both were
corrected to their \texttt{\_UNORM} equivalents.

\section{Buffers: VertexBuffer and IndexBuffer}

\cnaclass{VertexBuffer} and \cnaclass{IndexBuffer} both inherit \cnaclass{GraphicsResource}
and both maintain a private \texttt{cpuShadow\_} byte buffer purely to support
\texttt{GetData()} without requiring a real per-backend GPU readback path --- since nothing
in CNA's own pipeline ever writes GPU-side data back into one of these buffers, a CPU-side
shadow copy is a faithful, sufficient implementation of read-your-own-writes semantics.
\cnaclass{VertexBuffer::SetData} / \texttt{GetData} are typed per concrete vertex struct ---
\cnaclass{VertexPositionColor},
\cnaclass{VertexPositionColorTexture},
\cnaclass{VertexPositionNormalTexture}, and
\cnaclass{VertexPositionTexture} --- plus a
\texttt{NOXNA} skinned variant and a raw \texttt{SetDataRaw(data, count, stride)} escape hatch
for layouts with no dedicated XNA struct). \cnaclass{DynamicVertexBuffer} and
\cnaclass{DynamicIndexBuffer} add \texttt{IsContentLost} (always \texttt{false}) and
\texttt{SetData(..., SetDataOptions)} overloads --- but the \texttt{dynamic} construction flag
and the \texttt{Discard} / \texttt{NoOverwrite} streaming hints are both accepted purely for API
conformance and currently ignored by every backend: static and dynamic buffers take the
identical GPU upload path today, regardless of the hint given.

\section{Vertex layout: chosen by stride, not by declaration}
\label{sec:vertex-layout}

The single most important fact in this section is one that is easy to miss reading the API in
isolation: \textbf{every hardware backend selects its GPU vertex-attribute layout from the raw
byte stride of the bound vertex buffer, not from the \cnaclass{VertexDeclaration}'s own
element list.} \cnaclass{VertexDeclaration} is stored and used to compute the stride, and its
individual \cnaclass{VertexElementFormat} / \cnaclass{VertexElementUsage} mapping tables are
each individually correct on EasyGL, Vulkan, and BGFX --- but they are consulted only as an
audit reference today, not by the active draw dispatch. Exactly five hardcoded strides render
correctly: 16 bytes (\cnaclass{VertexPositionColor}), 20 (\cnaclass{VertexPositionTexture}),
24 (\cnaclass{VertexPositionColorTexture}), 32 (\cnaclass{VertexPositionNormalTexture}), and
52 (a custom skinned layout). Anything else falls back differently per backend: EasyGL
renders position-only with a logged warning; Vulkan compiles no pipeline at all and silently
skips the draw; BGFX falls back to position-plus-color-only, treating anything past that as
padding, which produces visibly wrong shading if the shader in question actually reads a
normal or UV attribute from that padding. SDL\_Renderer has no 3D vertex pipeline at all, so
\texttt{CreateVertexBuffer()} throws immediately, consistent with its documented 2D-only
scope.

If you are authoring a custom vertex layout for a ported game (this book's migration chapter
returns to this), the practical rule this section implies is blunt: stick to one of the five
recognized strides, or verify your chosen backend's fallback behavior directly, because the
declaration you wrote is not what determines what actually reaches the GPU.
